% Detailed solutions for VI/40 — Cartier Divisors.
% Canonical exercise statements remain in chapter.tex.

% VI40-DETAILED-EXERCISE-01
\paragraph{Solution VI.40.E01.}
Let \(D\) and \(E\) be represented on a common refinement by local equations \(f_i\) and \(g_i\).  Define \(D+E\) by \(f_i g_i\).  On an overlap,
\[
\frac{f_i g_i}{f_j g_j}=\frac{f_i}{f_j}\frac{g_i}{g_j}\in\OO_X^\times,
\]
so the product equations again define a Cartier divisor.  The zero divisor is represented by \(1\), and \(-D\) is represented by \(f_i^{-1}\), because
\[
\frac{f_i^{-1}}{f_j^{-1}}=\left(\frac{f_i}{f_j}\right)^{-1}
\]
is a unit.  Associativity and commutativity reduce to multiplication in \(\KX^\times\).  Thus \(\CaDiv(X)\) is an abelian group.

% VI40-DETAILED-EXERCISE-02
\paragraph{Solution VI.40.E02.}
If every local equation \(f_i\) belongs to \(\OO_X(U_i)^\times\), then its class in \(\KX^\times/\OO_X^\times\) is the identity class.  Hence all local representatives are zero in the quotient sheaf and glue to the zero global section.  Therefore the Cartier divisor is zero.

% VI40-DETAILED-EXERCISE-03
\paragraph{Solution VI.40.E03.}
On \(\AA^1_k\), the rational function
\[
f=t^3(t-1)^{-2}
\]
has order \(3\) at \(t=0\), order \(-2\) at \(t=1\), and order \(0\) at every other height-one prime.  Therefore the Cartier-to-Weil map sends the principal Cartier divisor \((f)_{\mathrm{Cart}}\) to
\[
3[0]-2[1].
\]
There is no contribution from infinity because the ambient scheme is affine.

% VI40-DETAILED-EXERCISE-04
\paragraph{Solution VI.40.E04.}
The coordinate ring \(k[x_1,\ldots,x_n]\) is a domain, so every nonzero element is a non-zero-divisor.  In particular \(x^m\) is a non-zero-divisor for \(m\ge1\).  Hence the principal ideal \((x^m)\) defines an effective Cartier divisor.  Its reduced support is \(V(x)\), while scheme-theoretically it is the \(m\)-fold thickening cut out by \((x^m)\).

% VI40-DETAILED-EXERCISE-05
\paragraph{Solution VI.40.E05.}
In
\[
A=k[x,y]/(xy),
\]
we have \(xy=0\) with both \(x\) and \(y\) nonzero.  Thus multiplication by \(x\) kills the nonzero class of \(y\), and multiplication by \(y\) kills the nonzero class of \(x\).  Hence neither is a non-zero-divisor.  Therefore neither \(V(x)\) nor \(V(y)\) is an effective Cartier divisor, despite each being principal as a closed subscheme.

% VI40-DETAILED-EXERCISE-06
\paragraph{Solution VI.40.E06.}
If \(D\) is effective Cartier with local equation \(f_i\), then \(f_i\) is a non-zero-divisor.  Every positive power \(f_i^n\) is again a non-zero-divisor: if \(f_i^n a=0\), repeated injectivity of multiplication by \(f_i\) gives \(a=0\).  The equations \(f_i^n\) represent \(nD\), so \(nD\) is effective.

% VI40-DETAILED-EXERCISE-07
\paragraph{Solution VI.40.E07.}
On an open set where \(D=V(f)\) with \(f\) a non-zero-divisor,
\[
\OO_X(-D)=f\OO_X
\]
by the divisor-sheaf construction.  But the ideal sheaf of the effective Cartier divisor is exactly \((f)\).  If \(f_i=u_{ij}f_j\) on an overlap with \(u_{ij}\) a unit, then \(f_i\OO=f_j\OO\), so the local identifications glue.  Hence
\[
\OO_X(-D)=\IIcal_D.
\]

% VI40-DETAILED-EXERCISE-08
\paragraph{Solution VI.40.E08.}
Let \(D\) and \(E\) have local equations \(f_i\) and \(g_i\).  Then \(\OO_X(D)\) and \(\OO_X(E)\) are locally generated inside \(\KX\) by \(f_i^{-1}\) and \(g_i^{-1}\).  Multiplication gives
\[
(f_i^{-1}\OO)\otimes(g_i^{-1}\OO)\longrightarrow(f_i g_i)^{-1}\OO,
\qquad
f_i^{-1}\otimes g_i^{-1}\longmapsto(f_i g_i)^{-1}.
\]
This is an isomorphism locally, and on overlaps it respects the unit transition factors.  Therefore
\[
\OO_X(D+E)\cong\OO_X(D)\otimes\OO_X(E).
\]

% VI40-DETAILED-EXERCISE-09
\paragraph{Solution VI.40.E09.}
The divisor \(V(t)\) on \(\AA^1_t\) has local equation \(t\).  Under \(g\), this equation becomes
\[
g^*t=u^4.
\]
Thus the pulled-back effective Cartier divisor is cut out by \((u^4)\), and therefore
\[
g^*V(t)=4V(u).
\]
The coefficient \(4\) is the ramification multiplicity of the map at \(u=0\).

% VI40-DETAILED-EXERCISE-10
\paragraph{Solution VI.40.E10.}
Let \(i:V(t)\hookrightarrow\AA^1_t\).  Pulling back the local equation \(t\) gives \(i^*t=0\).  A Cartier local equation must define an element of the unit group of the total quotient sheaf, whereas \(0\) is never invertible.  Hence \(i^*V(t)\) is not defined as a Cartier divisor.  Geometrically the entire source lies in the support of the divisor.

% VI40-DETAILED-EXERCISE-11
\paragraph{Solution VI.40.E11.}
Suppose \(D=(f)_{\mathrm{Cart}}\) is principal.  Whenever the pullback is defined, the single rational equation \(f\) pulls back to the rational unit \(g^*f\).  Therefore
\[
g^*D=(g^*f)_{\mathrm{Cart}},
\]
which is principal.  This also shows directly that pullback respects addition of principal divisors because pullback of rational functions respects multiplication.

% VI40-DETAILED-EXERCISE-12
\paragraph{Solution VI.40.E12.}
Near the generic point \(\eta_Z\), two equivalent local equations satisfy \(f'=uf\) for some \(u\in\OO_{X,\eta_Z}^\times\).  Since \(\OO_{X,\eta_Z}\) is a DVR, every unit has valuation zero.  Thus
\[
\operatorname{ord}_Z(f')
=\operatorname{ord}_Z(u)+\operatorname{ord}_Z(f)
=\operatorname{ord}_Z(f).
\]
Hence the Cartier-to-Weil coefficient is independent of the chosen local equation.

% VI40-DETAILED-EXERCISE-13
\paragraph{Solution VI.40.E13.}
If \(D\) is effective Cartier, its local equation near \(\eta_Z\) is a regular non-zero-divisor \(f\).  In the codimension-one DVR \(\OO_{X,\eta_Z}\), every regular element has nonnegative valuation.  Hence every coefficient
\[
n_Z(D)=\operatorname{ord}_Z(f)
\]
is nonnegative.  Therefore the associated Weil divisor is effective.

% VI40-DETAILED-EXERCISE-14
\paragraph{Solution VI.40.E14.}
For a principal Cartier divisor \((f)_{\mathrm{Cart}}\), the local equation at every codimension-one point is the same rational function \(f\).  The Cartier-to-Weil coefficient along \(Z\) is therefore \(\operatorname{ord}_Z(f)\).  Consequently
\[
[(f)_{\mathrm{Cart}}]_{\mathrm W}
=\sum_Z\operatorname{ord}_Z(f)[Z]
=\operatorname{div}_X(f).
\]

% VI40-DETAILED-EXERCISE-15
\paragraph{Solution VI.40.E15.}
The ring \(k[x_1,\ldots,x_n]\) is a UFD.  Every height-one prime ideal is therefore generated by an irreducible polynomial \(p\).  The corresponding prime Weil divisor \(V(p)\) is locally, in fact globally, cut out by the non-zero-divisor \(p\), so it is Cartier.  Since a Weil divisor is a finite integer combination of prime divisors, every Weil divisor on \(\AA^n_k\) is Cartier.

% VI40-DETAILED-EXERCISE-16
\paragraph{Solution VI.40.E16.}
A Noetherian regular local ring is a UFD.  Thus a regular Noetherian integral scheme is locally factorial.  Every height-one prime ideal in each local ring is principal, so every prime Weil divisor is locally principal by a non-zero-divisor.  Hence every prime Weil divisor is Cartier, and therefore every finite integer combination of prime Weil divisors is Cartier.

% VI40-DETAILED-EXERCISE-17
\paragraph{Solution VI.40.E17.}
For \(H=V_+(x_0)\), use
\[
f_i=\frac{x_0}{x_i}\quad\text{on }U_i=D_+(x_i).
\]
On \(U_i\cap U_j\),
\[
\frac{f_i}{f_j}
=\frac{x_0/x_i}{x_0/x_j}
=\frac{x_j}{x_i}.
\]
Both \(x_i\) and \(x_j\) are nonzero on the overlap, so \(x_j/x_i\) is a regular unit there.  Thus the equations satisfy the Cartier compatibility condition.

% VI40-DETAILED-EXERCISE-18
\paragraph{Solution VI.40.E18.}
For the hyperplane \(H=V_+(x_0)\), the local Cartier equations are \(f_i=x_0/x_i\).  Hence \(\OO(H)\) is locally generated by
\[
e_i=f_i^{-1}=\frac{x_i}{x_0}.
\]
On \(U_i\cap U_j\),
\[
e_i=\frac{x_i}{x_j}e_j.
\]
These are precisely the standard transition functions for \(\OO_{\PP^n}(1)\).  Therefore
\[
\OO_{\PP^n}(H)\cong\OO_{\PP^n}(1).
\]

% VI40-DETAILED-EXERCISE-19
\paragraph{Solution VI.40.E19.}
If \(D=(f)_{\mathrm{Cart}}\), then
\[
\OO_X(D)=f^{-1}\OO_X\subseteq\KX.
\]
Multiplication by \(f\) defines an \(\OO_X\)-linear map
\[
f^{-1}\OO_X\longrightarrow\OO_X
\]
with inverse multiplication by \(f^{-1}\).  Thus \(\OO_X(D)\cong\OO_X\).  Principal Cartier divisors therefore yield trivial divisor line bundles.

% VI40-DETAILED-EXERCISE-20
\paragraph{Solution VI.40.E20.}
If \(D-E=(f)_{\mathrm{Cart}}\), then by the tensor law,
\[
\OO_X(D)
\cong \OO_X(E)\otimes\OO_X(D-E)
\cong \OO_X(E)\otimes\OO_X((f)).
\]
The last factor is trivial by the preceding exercise, so
\[
\OO_X(D)\cong\OO_X(E).
\]
The isomorphism is concretely multiplication by a suitable rational function corresponding to \(f\).

% VI40-DETAILED-EXERCISE-21
\paragraph{Solution VI.40.E21.}
Let \(D\) be effective with local equation \(f_i\).  In \(\OO_X(D)|_{U_i}=f_i^{-1}\OO_{U_i}\), use the local generator \(e_i=f_i^{-1}\).  The canonical rational section \(s_D=1\) satisfies
\[
1=f_i e_i.
\]
Because \(f_i\in\OO_X(U_i)\), the coefficient of \(s_D\) in this trivialization is regular.  Hence \(s_D\) is a regular global section.

% VI40-DETAILED-EXERCISE-22
\paragraph{Solution VI.40.E22.}
In the same trivialization \(e_i=f_i^{-1}\), the canonical section is represented by the regular function \(f_i\).  Therefore its zero subscheme on \(U_i\) is cut out by \((f_i)\).  Those ideals are exactly the local defining ideals of the effective Cartier divisor \(D\), and they glue.  Hence the zero subscheme of the canonical section is precisely \(D\).

% VI40-DETAILED-EXERCISE-23
\paragraph{Solution VI.40.E23.}
Choose local generators \(e_i\) of \(\mathcal L\) and write the rational section as \(s=f_i e_i\).  If \(f\in K(X)^\times\), then
\[
fs=(ff_i)e_i.
\]
Thus all Cartier local equations are multiplied by the same rational function \(f\).  Addition of Cartier divisors is multiplication of local equations, so
\[
\operatorname{div}(fs)=\operatorname{div}(s)+(f)_{\mathrm{Cart}}.
\]

% VI40-DETAILED-EXERCISE-24
\paragraph{Solution VI.40.E24.}
Let \(C\) be a smooth integral curve and \(P\in C\) a closed point.  Smoothness implies that \(\OO_{C,P}\) is a one-dimensional regular local ring, hence a DVR.  Its maximal ideal is therefore generated by a uniformizer \(\pi\), and \(\pi\) is a non-zero-divisor because the local ring is a domain.  On a sufficiently small neighborhood of \(P\), \(\pi\) cuts out \(P\) scheme-theoretically.  Thus \(P\) is an effective Cartier divisor.
